\section{Cross-Technique Comparison}
\label{sec:comparison}

% Target: ~1 page. This is the section that answers ``so what'' for
% the whole study; spend disproportionate effort here.

\subsection{Summary Table}
\label{sec:comp:table}
\Cref{tab:comparison} aggregates the best configuration from each
preceding section into one table. We deliberately use the
\emph{best} setting per family (not the average) because the question
is which family is worth investing in, not how robust each family is
to its own hyperparameter.

% Cross-technique summary table. Each row = best configuration of one
% regularization family. Replace -- with real numbers; bold winners.

\begin{table}[t]
  \centering
  \small
  \caption{Best configuration per regularization family on
  CIFAR-100 with ViT-Tiny/16. ``Ease'' is a qualitative judgment of
  implementation effort; ``Cost'' is wall-clock training time
  relative to baseline. Best per column in bold.
  \draftnote{Fill in once all sections have results. Bold winners.}}
  \label{tab:comparison}
  \begin{tabular}{lcccll}
    \toprule
    Regularization (best config)
      & Test acc (\%) $\uparrow$
      & Best val acc (\%) $\uparrow$
      & Gap reduction vs baseline (pp) $\uparrow$
      & Ease
      & Train cost \\
    \midrule
    None (baseline)            & -- & -- & 0.0 & --  & $1.0\times$ \\
    Weight decay (best $\lambda$)
                               & -- & -- & --  & easy   & $\approx 1.0\times$ \\
    Dropout (best rate)        & -- & -- & --  & easy   & $\approx 1.0\times$ \\
    Augmentation (best combo)  & -- & -- & --  & medium & -- \\
    Label smoothing (best $\epsilon$)
                               & -- & -- & --  & easy   & $\approx 1.0\times$ \\
    Early stopping (best patience)
                               & -- & -- & --  & easy   & $< 1.0\times$ \\
    \bottomrule
  \end{tabular}
\end{table}


\subsection{Visual Summary}
\label{sec:comp:bar}
\begin{figure}[t]
  \centering
  % \includegraphics[width=0.7\linewidth]{comparison_bar.png}  % uncomment when ready
  \fbox{\rule{0pt}{4cm}\rule{0.65\linewidth}{0pt}}
  \caption{Best test accuracy per regularization family, with the
  no-regularization baseline as reference. \draftnote{Drop
  \texttt{figures/comparison\_bar.png} and swap in the commented
  \texttt{\textbackslash includegraphics} above. Sort bars by test
  accuracy and annotate each with the absolute improvement over
  baseline.}}
  \label{fig:comparison-bar}
\end{figure}

\subsection{Which Family Matters Most for ViT?}
\label{sec:comp:discussion}
The four regularizer families correspond to four levels at which a
model can be constrained:
\begin{itemize}
  \item \textbf{Data-level}: augmentation (\cref{sec:augmentation}).
  \item \textbf{Parameter-level}: weight decay (\cref{sec:wd}),
        dropout (\cref{sec:dropout}).
  \item \textbf{Output-level}: label smoothing (\cref{sec:lsearly}).
  \item \textbf{Training-process-level}: early stopping
        (\cref{sec:lsearly}).
\end{itemize}
\draftnote{Write 1--2 paragraphs that argue from the numbers in
\cref{tab:comparison} which level pays off most for a small ViT in
the small-data fine-tuning regime. Be specific: cite percentage-point
improvements rather than vague language. If two families tie, say so
and discuss the tradeoff (e.g. data augmentation costs more compute
per epoch; weight decay is free).}

\subsection{Stability}
\label{sec:comp:stability}
\draftnote{If multi-seed runs are available, report the standard
deviation of test accuracy per family and discuss which is most
stable. Otherwise, explicitly note that all numbers are single-seed
and that the ranking should be interpreted accordingly. This should be
stated consistently with the limitations discussed in the per-method
results sections.}
